\section{魏：最先跑起来的中原国家，最先感到疲惫}

从安邑向东望去，黄河与汾水之间的道路把魏推到列国往来的中央。三家灭智、周室册命之后，魏接过的不只是旧晋的一部分城邑和田土，也接过了一个不再容许卿族各自经营的中原。西面隔河是秦，东面通齐，南北又与韩、赵、楚相接；粮食、士人和兵车能够沿着这些道路汇集，敌军也同样能够沿着它们逼近。魏文侯一朝之所以在后世叙事中显得格外明亮，首先不在于某一位臣子忽然发明了新制度，而在于这个新国家较早地把田亩、法令、官吏和兵员视为同一场竞争里的资源。它以\eventref{WQ-0403-THREEJIN}{战国·三家分晋}所确认的名分为起点，却必须靠实际的征发与组织来守住这份名分。

\begin{event}{前408年前后}{李悝传统与魏国的早期整顿}{年代与细节有争议}{魏}\eventanchor{WQ-0408-WEI}{战国·魏国早期改革}
《史记·魏世家》把李悝列在魏文侯所任用的人物之中；较晚的《汉书·食货志》则把“尽地力”与平籴的议论系于他。两种叙事共同保存了一项重要的历史记忆：当谷物丰歉足以摇动价格、徭役和军粮时，国家不能只在收成之后向封邑索取，而要设法看见土地的产出、储藏与流通。魏的优势正在这里逐步形成：旧贵族的门第仍有分量，但能被官府统计、调度和供给的资源，开始比家族私下的从属关系更能决定一国的行动速度。

不过，后人常把这段历史收束成“李悝作《法经》、魏国遂变法”的整齐开端，材料并不容许这样的肯定。今见《法经》篇目主要出自后出的法制史追述，原书早已不存；“平籴”的办法也不能据此倒推出一套已可逐条复原的魏国法规。从现存材料看，更稳妥的说法是：魏国早期确有把农政、刑名和任官纳入国家竞争的趋势，李悝成为这一趋势最有力的后世符号。个人的履历、政策的形成和制度的长期运行，不能被压成同一件事。
\end{event}

\begin{debate}[李悝与《法经》：传统可以指路，不能代替原件]
《史记》《汉书·食货志》提供了李悝与魏文侯、农政议论相连的叙述；《晋书·刑法志》等后出材料保存了《法经》六篇的法制传统。它们足以说明魏被记忆为早期制度竞争的实验场，却不足以确证每一条法文、每一次颁行和全部政策都出自李悝一人。出土秦律能够帮助理解先秦法律文书化的深度，但不是魏国改革的直接档案。
\end{debate}

李悝传统所指向的，是把散在乡里的收成与劳力纳入可计算秩序的努力；吴起的故事所指向的，则是把同一批可征发的人变成可反复使用的军队。两者未必是一个人设计的一套蓝图，却在魏的中原处境中彼此需要：没有较稳定的供给，精锐无从久驻；没有能在边境取胜的军队，田地与市场也难以免于他国征取。

\begin{event}{魏文侯、武侯时期}{吴起与魏武卒}{年代与细节有争议}{魏}\eventanchor{WQ-WEI-WUQI}{战国·魏武卒}
《史记·孙子吴起列传》说吴起在魏用兵有功，后世遂以“魏武卒”概括魏军的锐利。《吴子》又描绘负甲、持兵、行军的选拔条件，仿佛一张可以逐项核验的征募清单。可这部兵书的成书层累和这些标准的适用年代都存在疑问，不能把它直接当作魏文侯、武侯朝的一份军籍簿。可以把握的，是魏确实在战国早期以更严密的选拔、训练、赏罚和供给来提高军队的可用性；至于“武卒”的人数、装备和全部考核办法，仍须留在传说与制度史之间审慎辨认。

这里的人事转折并不神秘。像吴起这样的外来行动者能够在魏获得位置，反映的是君主需要能够组织战争的专业才能；而这类才能一旦与旧有贵族、君权和军功分配发生摩擦，也不可能只凭战场声望安然存在。史书随后写吴起离魏入楚，所保留的正是这种政治处境的结果，而非一段可以补写出心理独白的个人传奇。
\end{event}

魏军的锋芒使它能够向外施压，也使中原诸国更早学会以联盟回应它。前354年，魏军围赵都邯郸，赵向齐求援。齐不必在赵城下同魏军硬碰，便可选择把战事引回魏所倚重的腹地；《史记·孙子吴起列传》所叙田忌、孙膑的行动，由此成为“围魏救赵”的经典模型。战场在桂陵，胜负所打断的却是魏以一场围城迫使赵屈服的盘算。

\begin{event}{前354年起}{邯郸危局与桂陵受挫}{可靠纪年}{魏、赵、齐}\eventanchor{WQ-0354-GUILIN}{战国·桂陵之战}
桂陵的故事之所以传得最广，正在于它把复杂的战略选择压缩为一条迂回路线。但围邯郸、齐军出动与桂陵交锋在不同史书中的编年并非完全整齐，后出的兵家叙事又强化了孙膑运筹的戏剧性。能够确定的是，齐的介入迫使魏在赵地之外顾及本土和补给线；\eventref{WQ-0354-GUILIN}{战国·桂陵之战}的直接后果，不是魏立刻失去一切，而是它再也不能把中原的每一场战争当成两国之间的单线较量。
\end{event}

十余年后，魏攻韩，齐又以救援者的身份出兵。马陵一役在前342年的传统编年中成为魏军受创最深的节点。传世叙事写孙膑减灶诱敌、庞涓自尽，画面锐利，却主要依赖《史记》与《战国策》的后出铺陈；银雀山汉简所见《孙膑兵法》证明相关兵学传统并非全然无根，却不能逐句为这场战役作证。应当保留的，不是一则奇谋的细节，而是魏军在长期外战、强邻合纵与补给压力交织之下，已难维持早期那种由它选择战场的主动。

\begin{event}{前342年}{马陵之战与魏的战略收缩}{可靠纪年}{魏、韩、齐}\eventanchor{WQ-0342-MALING}{战国·马陵之战}
马陵之后，魏仍是大国，后来还与齐互称王；败战并不等于国家骤然坍塌。它改变的是力量的相对位置：齐证明自己能够越出东方干预中原，韩、赵也不再只是魏的附庸式邻国，西面的秦则在河西与关中的竞争中等待机会。魏的地理曾使它最先聚拢道路和资源，如今也使它必须在更多方向上分兵、结交和付出财政成本。战国的中心没有消失，只是不再由魏独占。
\end{event}

\begin{sourcenote}[桂陵、马陵与魏国衰势]
两役的基本对抗及其前后政治影响，可据《史记·孙子吴起列传》《魏世家》与《战国策》互相参照；人物谋略、伤亡数字和具体战术则带有显著的后世叙事加工。现代研究通常把目光放在齐魏竞争、行军与后勤条件、孙膑故事的成形过程，而不把“减灶”一类细节当作无需检验的战场实录。
\end{sourcenote}

\begin{turningpoint}[制度得失：先行者为何也会落后]
魏较早把农、法、兵连成国家能力，因此取得了先发的速度；承担代价的却是被征发的乡里、持续供给军队的财政，以及必须面对多面威胁的中原社会。桂陵与马陵不是制度失效的简单判决，而是地缘、联盟和战争消耗对先发优势的反馈。魏未能把早期的组织能力转化为不受制衡的霸权，后来秦所继承并重组的，正是这种竞争留下的难题。
\end{turningpoint}
